\documentclass[11pt,a4paper]{article}
\usepackage[margin=1in]{geometry}
\usepackage{amsmath,amssymb,amsfonts}
\usepackage{booktabs}
\usepackage{hyperref}
\usepackage{xcolor}

\hypersetup{colorlinks=true, linkcolor=blue, citecolor=blue, urlcolor=blue}

\title{Response to Collaborator Questions\\on the APN Report}
\author{Zhixin Lu}
\date{April 2026}

\begin{document}
\maketitle

% ============================================================================
\section{Question 1: Delta Rule vs.\ Oja's Rule}

\textbf{Question:} ``Your correction term in Eq.\ (1) looks kind of like Oja's rule, but `opposite': in your notation, Oja's normalizing term would be proportional to $v^2$ (or perhaps $o^2$). What's the motivation behind your correction term? Would switching to an Oja-like rule not be advantageous? (See \href{https://arxiv.org/abs/2408.08408v3}{arXiv:2408.08408v3}.)''

\subsection{The ``$W + M$'' Story: APN as Plastic Weight Augmentation}

The core idea of APN is simple: each layer has a static weight $W$ (trained by backprop) augmented by a plastic state $S$ (updated online during the forward pass).
The layer output is:
\begin{equation}\label{eq:output}
\mathbf{o}_t = \underbrace{W^\top \text{act}(\mathbf{x}_t)}_{\text{static: } \mathbf{v}_t} + \underbrace{S_t^\top \text{act}(\mathbf{x}_t)}_{\text{plastic recall}} = (W + S_t)^\top \text{act}(\mathbf{x}_t)
\end{equation}
The effective weight at time $t$ is $W + S_t$, where $S_t$ accumulates experience from the sequence seen so far.
$W$ is fixed within a forward pass (updated only by the training loss); $S_t$ adapts \emph{within} the sequence via a Hebbian-family learning rule.

The plastic state decays and updates as:
\begin{equation}\label{eq:update}
S_t = \lambda \, S_{t-1} + \beta \, \mathbf{k}_t \otimes \mathbf{h}_t, \qquad \beta = \frac{\eta(1-\lambda)}{d}
\end{equation}
where $\mathbf{k}_t = \text{act}(\mathbf{x}_t)$ is the key (pre-synaptic activity) and $\mathbf{h}_t$ is the post-synaptic learning signal.
The entire design reduces to one question: \textbf{what is $\mathbf{h}_t$?}

\subsection{Original APAN Intuition vs.\ Implemented APN}

\subsubsection{Original APAN: Hebbian on the full output (unstable)}

The original neuroscience intuition was to use the layer's own output as the learning signal:
\begin{equation}
\mathbf{h}_t^{\text{APAN}} = \mathbf{o}_t = \mathbf{v}_t + S_{t-1}^\top \mathbf{k}_t \qquad \text{(Hebbian: pre $\times$ post)}
\end{equation}
This gives a pure Hebbian rule $\Delta S = \beta \, \mathbf{k} \otimes \mathbf{o}$---the connection strengthens when pre and post are coactive.
Substituting $\mathbf{o} = \mathbf{v} + S^\top \mathbf{k}$:
\begin{equation}
S_t = \lambda \, S_{t-1} + \beta \, \mathbf{k}_t \otimes (\mathbf{v}_t + S_{t-1}^\top \mathbf{k}_t)
\end{equation}

\textbf{Problem: positive feedback.}
Project the recurrence onto a repeated key direction, letting $r_t = S_t^\top \mathbf{k}$ denote the recall for that key:
\begin{equation}
r_t = \underbrace{(\lambda + \beta \|\mathbf{k}\|^2)}_{\alpha} \, r_{t-1} + \beta \|\mathbf{k}\|^2 \, v_{\text{proj}}
\end{equation}
Stability requires $|\alpha| < 1$.
With typical values ($\lambda=0.99$, $\eta=1$, $d=100$, $\|\mathbf{k}\|^2 \approx d = 100$ for unit-variance activations):
\[
\alpha = 0.99 + \frac{1.0 \times 0.01}{100} \times 100 = 0.99 + 0.01 = 1.00
\]
The system is \textbf{exactly at the stability boundary}.
For keys with $\|\mathbf{k}\|^2 > d$ (common without key normalization), $\alpha > 1$ and the state diverges exponentially.
The mechanism: $S^\top \mathbf{k}$ is large $\to$ output $\mathbf{o}$ is large $\to$ we store more into $S$ $\to$ next recall is larger $\to$ runaway.

\subsubsection{Implemented APN: Hebbian on the prediction error (stable)}

The implemented APN uses the \textbf{prediction error} as the learning signal:
\begin{equation}
\mathbf{h}_t^{\text{APN}} = \mathbf{v}_t - S_{t-1}^\top \mathbf{k}_t \qquad \text{(target $-$ recall = surprise)}
\end{equation}
giving:
\begin{equation}\label{eq:apn-implemented}
S_t = \lambda \, S_{t-1} + \beta \, \mathbf{k}_t \otimes (\mathbf{v}_t - S_{t-1}^\top \mathbf{k}_t)
\end{equation}

The same projection onto a repeated key gives:
\begin{equation}
r_t = \underbrace{(\lambda - \beta \|\mathbf{k}\|^2)}_{\alpha} \, r_{t-1} + \beta \|\mathbf{k}\|^2 \, v_{\text{proj}}
\end{equation}
Now $\alpha = 0.99 - 0.01 = 0.98 < 1$ --- \textbf{unconditionally stable}.
The sign flip provides \emph{negative feedback}: larger recall $\to$ smaller error $\to$ smaller update $\to$ convergence.
At the fixed point $S^\top \mathbf{k} = \mathbf{v}$, the error is zero and the memory stops updating---it has learned the association perfectly.

\subsection{Is This Still Hebbian?}

Yes. The update $\Delta S = \beta \, \mathbf{k} \otimes \mathbf{h}$ is Hebbian in the broad sense: a rank-one outer product of pre-synaptic activity ($\mathbf{k}$) and post-synaptic activity ($\mathbf{h}$).
The post-synaptic signal $\mathbf{h} = \mathbf{v} - S^\top\mathbf{k}$ is simply a \emph{different neuron's output}---specifically, a \textbf{prediction error neuron} that computes ``what the static pathway says'' minus ``what the plastic memory already recalls.''

This is biologically plausible and well-precedented:
\begin{itemize}
\item \textbf{Predictive coding:} Synapses update proportional to prediction errors, not raw activity. Cortical microcircuits are hypothesized to separate representation neurons from error neurons.
\item \textbf{Cerebellar learning:} Parallel fiber$\to$Purkinje cell synapses update via climbing-fiber error signals, not via Purkinje output.
\item \textbf{Decorrelation / anti-Hebbian circuits:} Lateral inhibition networks that prevent redundant storage.
\end{itemize}

The ``$W + M$'' intuition is fully preserved:
\begin{itemize}
\item Effective weight at time $t$: $W + S_t$ \checkmark
\item Output: $(W + S_t)^\top \text{act}(\mathbf{x})$ \checkmark
\item Plasticity: local, activity-dependent, rank-one updates \checkmark
\item The only refinement: the post-synaptic learning signal is the \emph{surprise} (novelty) rather than the \emph{total output}.
Expected inputs produce zero update; surprising inputs drive learning.
\end{itemize}

\subsection{How Does This Relate to DeltaNet?}

DeltaNet (Yang et al., 2024) uses the same mathematical recurrence:
\begin{equation}
S_t = S_{t-1} + \beta_t \, \mathbf{k}_t \otimes (\mathbf{v}_t - S_{t-1}^\top \mathbf{k}_t)
\end{equation}
The APN differs from DeltaNet in three ways:

\begin{table}[h]
\centering
\caption{APN vs.\ DeltaNet: same kernel, different model.}
\begin{tabular}{lll}
\toprule
& \textbf{DeltaNet} & \textbf{APN (ours)} \\
\midrule
Gating & No decay ($\lambda = 1$) & Learnable decay $\lambda \in (0,1)$ \\
 & or separate gate function & via sigmoid parameterization \\
\midrule
Learning rate $\beta$ & $\beta_t = \sigma(\mathbf{w}^\top \mathbf{x}_t)$ & $\beta = \eta(1-\lambda)/d$ \\
 & input-dependent, per-token & scalar, shared across tokens \\
\midrule
Keys / Queries & Separate $\mathbf{q} = W_q\mathbf{x}$, $\mathbf{k} = W_k\mathbf{x}$ & Tied: $\mathbf{q} = \mathbf{k} = \text{act}(\mathbf{x})$ \\
 & linear projections & optional nonlinear activation \\
\midrule
Values & $\mathbf{v} = W_v\mathbf{x}$ (separate proj.) & $\mathbf{v} = W\, \text{act}(\mathbf{x})$ (same as static) \\
 & standard attention-style & ``$W + M$'' interpretation \\
\midrule
Motivation & Efficient attention with & Synaptic plasticity: $W$ is \\
 & better memory capacity & the genome, $S$ is experience \\
\bottomrule
\end{tabular}
\end{table}

\textbf{Key insight:} APN reinterprets the delta-rule recurrence through a neuroscience lens.
DeltaNet was designed as an efficient attention replacement (keys $\neq$ queries, input-dependent $\beta$, no decay).
APN constrains the architecture to be biologically interpretable: tied keys/queries (same pre-synaptic signal for storage and retrieval), scalar learning rate (uniform plasticity), and explicit decay (forgetting / memory lifetime control).
These constraints happen to make APN simpler (fewer parameters) while achieving competitive performance.

The shared kernel (\texttt{chunk\_gated\_delta\_rule}) is a fortunate alignment: the delta-rule recurrence that provides stability in the APN's ``$W + M$'' picture is exactly the recurrence that DeltaNet uses for efficient attention.
APN adds the gated decay ($\lambda$-based exponential forgetting) which DeltaNet in its original form does not have, but which the \emph{gated} delta-rule kernel supports.

\subsection{Why Not Oja's Rule?}

Oja's rule ($\Delta S \propto \mathbf{k}\mathbf{o}^\top - \|\mathbf{o}\|^2 S$) performs online PCA---it extracts principal components while keeping weights bounded.
This is a different objective:
\begin{itemize}
\item \textbf{Oja's rule} $\to$ dimensionality reduction / feature extraction.
\item \textbf{Delta rule} $\to$ associative recall (content-addressable memory).
\end{itemize}
APN needs memory---it must store and retrieve specific key-value associations from the sequence.
Oja's rule would find dominant directions of variation in the keys, but it cannot retrieve a specific value given a specific key.

That said, the preprint (\href{https://arxiv.org/abs/2408.08408v3}{arXiv:2408.08408v3}) proposes Oja-like normalization for stabilizing hidden representations.
This is \emph{complementary} to APN: one could apply Oja's normalization to the representations $\mathbf{x}$ \emph{before} they enter the delta-rule memory, combining feature normalization (Oja) with associative storage (delta rule).
This is an unexplored direction worth investigating.

% ============================================================================
\section{Question 2: Why Not Use $|\eta|$?}

\textbf{Question:} ``In any case, why not having $|\eta|$ in front of the correction term? That way you could still have anti-Hebbian when needed.''

\subsection{This Is a Good Suggestion}

Using $|\eta|$ (or equivalently, parameterizing as $\eta = |\hat{\eta}|$ for a free parameter $\hat{\eta}$) would:
\begin{enumerate}
\item \textbf{Prevent the multi-head divergence} we observed, where $\eta$ drifting negative caused the correction term to reverse sign, creating a positive feedback loop.
\item \textbf{Still allow anti-Hebbian corrections when needed}, because the sign of the full update depends on the sign of $(\mathbf{v}_t - S_{t-1}^\top \mathbf{k}_t)$---when the recall $S^\top \mathbf{k}$ overshoots the target $\mathbf{v}$, the correction naturally becomes anti-Hebbian (pushing the memory in the opposite direction). So $|\eta| > 0$ doesn't prevent anti-Hebbian corrections; it prevents the \emph{learning rate itself} from flipping sign.
\end{enumerate}

\subsection{Why We Didn't Use $|\eta|$ in Current Experiments}

The original motivation was to keep the parameterization as unconstrained as possible to observe what the model learns freely.
For single-head APN ($H=1$, $D=100$), the unconstrained $\eta$ works fine---it stays near its initialization ($\eta = 1.0$) and never goes negative, because the gradients are well-behaved at that scale (see Figure~7 in the report, where $\eta$ barely moves from 1.0 over 200 epochs).
The instability only manifests in the multi-head case ($H=4$, per-head $d=25$) where the small state dimension amplifies noise.

\subsection{Comparison of Constraint Options}

\begin{table}[h]
\centering
\caption{Parameterization options for $\eta$.}
\begin{tabular}{lcccl}
\toprule
\textbf{Parameterization} & $\eta > 0$? & \textbf{Smooth?} & \textbf{Bounded?} & \textbf{Notes} \\
\midrule
$\eta$ (unconstrained) & No & Yes & No & Diverges for multi-head \\
$|\eta|$ & Yes & No (at 0) & No & Simple; non-smooth at boundary \\
$\text{softplus}(\hat{\eta}) = \log(1+e^{\hat{\eta}})$ & Yes & Yes & No & Preferred; $\approx \hat{\eta}$ for $\hat{\eta} \gg 0$ \\
$\sigma(\hat{\eta})$ (sigmoid) & Yes & Yes & $(0,1)$ & Too restrictive; can't exceed 1 \\
$e^{\hat{\eta}}$ & Yes & Yes & No & Risk of exponential blowup \\
\bottomrule
\end{tabular}
\end{table}

Our recommendation going forward is \textbf{softplus} rather than $|\eta|$:
\begin{itemize}
\item Smooth everywhere (no gradient discontinuity at $\eta = 0$).
\item Approximately linear for $\eta \gg 0$: doesn't distort dynamics when $\eta$ is in a reasonable range.
\item Naturally prevents $\eta$ from reaching exactly 0 (which would kill plasticity entirely).
\end{itemize}

That said, $|\eta|$ is a perfectly valid choice and would likely fix the multi-head instability.
The gradient discontinuity at 0 is rarely a practical problem if $\eta$ is initialized away from 0 (e.g., at 1.0).

\textbf{A related note:} For $\lambda$ (the decay), we already use a sigmoid parameterization ($\lambda = \sigma(\hat{\lambda})$) to ensure $\lambda \in (0,1)$.
The asymmetry---constraining $\lambda$ but leaving $\eta$ unconstrained---was an oversight in our initial experiments.
Future runs will use softplus for $\eta$.

% ============================================================================
\section{Question 3: Pre-Norm vs.\ Post-Norm in APN}

\textbf{Question:} ``Is there any reason why you chose to use post-norm in your models? In Transformers pre-norm seems to be better (and you do use it, so it seems).''

\subsection{Clarification: APN Uses Branch-Norm, Not Standard Post-Norm}

Let me carefully describe the exact normalization used in each model as implemented in our code.

\subsubsection{Transformer (standard pre-norm)}

Our Transformer uses the standard pre-norm formulation:
\begin{align}
\mathbf{h} &\leftarrow \mathbf{h} + \text{MHA}(\text{LN}_1(\mathbf{h})) \\
\mathbf{h} &\leftarrow \mathbf{h} + \text{FFN}(\text{LN}_2(\mathbf{h}))
\end{align}
This is standard practice. LayerNorm is applied to the residual stream \emph{before} each sub-layer. The \texttt{TransformerLayer.forward()} returns the full residual delta (attention + FFN), and the outer loop adds it: \texttt{h = h + layer(h)}.

\subsubsection{APN layer (branch-norm)}

The APN layer applies LayerNorm \emph{to its output branch before the residual addition}:
\begin{equation}
\texttt{APNLayer}(\mathbf{h}) = \text{LN}\bigl(\underbrace{W \cdot \text{act}(\mathbf{h})}_{\text{static: } \mathbf{v}} + \underbrace{\text{delta\_recall}(\text{act}(\mathbf{h}))}_{\text{dynamic: } S^\top\mathbf{k}}\bigr)
\end{equation}
The residual is then added by the outer loop:
\begin{equation}
\mathbf{h} \leftarrow \mathbf{h} + \text{LN}\bigl(W \cdot \text{act}(\mathbf{h}) + S^\top \text{act}(\mathbf{h})\bigr)
\end{equation}

In code:
\begin{verbatim}
# Inside APNLayer.forward():
return self.o_norm(static_out + o_dyn.reshape(B, L, D))

# Inside SeqModel.forward():
h = h + layer(h)   # residual added externally
\end{verbatim}

This is \textbf{neither standard pre-norm} (\texttt{h + Layer(LN(h))}) \textbf{nor standard post-norm} (\texttt{LN(h + Layer(h))}).
It is sometimes called ``sub-LayerNorm'' or ``branch normalization'': the layer's output is normalized before being added to the residual stream, but the layer's \emph{input} is NOT normalized.

\subsubsection{APN + FFN (combined)}

When FFN blocks are added, the FFN itself uses pre-norm:
\begin{align}
\mathbf{h} &\leftarrow \mathbf{h} + \text{LN}_{\text{branch}}\bigl(\text{APN}(\mathbf{h})\bigr) & \text{(APN with branch-norm)} \\
\mathbf{h} &\leftarrow \mathbf{h} + \text{FFN}\bigl(\text{LN}_{\text{pre}}(\mathbf{h})\bigr) & \text{(FFN with pre-norm)}
\end{align}

In code:
\begin{verbatim}
# Inside SeqModel.forward() with use_ffn=True:
for layer, ffn in zip(self.layers, self.ffns):
    h = h + layer(h)    # APN output is branch-normed internally
    h = h + ffn(h)      # FFN has pre-norm internally: ffn(x) = MLP(LN(x))
\end{verbatim}

\subsubsection{Final norm}

After all layers, a final LayerNorm is applied to the last token before classification:
\begin{equation}
\text{logits} = W_{\text{cls}} \cdot \text{LN}_{\text{final}}(\mathbf{h}_{T})
\end{equation}

\subsection{Why Branch-Norm Instead of Pre-Norm for APN?}

\begin{enumerate}
\item \textbf{The key representation should not be normalized.}
In APN, the keys $\mathbf{k}_t = \text{act}(\mathbf{x}_t)$ enter both the memory query and the memory update.
If we applied pre-norm (LayerNorm on the input $\mathbf{h}$ before the layer), the keys would be normalized to unit variance at every layer.
This would compress the key space and interfere with the delta-rule dynamics---the memory needs to distinguish between inputs of different magnitudes to store distinct associations.
Our experimental results support this: \texttt{act=none} (identity, preserving full dynamic range) outperforms \texttt{tanh} (bounded keys) by 5--6 points for exactly this reason.
Pre-norm would have a similar compressive effect.

\item \textbf{Residual stream magnitude control.}
Branch-norm ensures that regardless of what happens inside the delta-rule computation (which involves a $D \times D$ state matrix that could amplify signals), the contribution added to the residual stream has bounded magnitude (zero mean, unit variance).
This provides training stability without interfering with the internal memory dynamics.

\item \textbf{The Transformer is different.}
Pre-norm works well for Transformers because the attention mechanism is inherently scale-sensitive (softmax is sensitive to logit magnitudes, and Q/K scaling is built in).
LayerNorm before attention ensures that the query/key dot products are well-conditioned.
For APN, the ``attention'' is implicit in the delta-rule recall $S^\top \mathbf{k}$, and the scaling is handled explicitly by the factor $\eta(1-\lambda)/d$---so input normalization is less critical.

\item \textbf{Depth considerations.}
Pre-norm's main advantage in Transformers is enabling stable training for very deep networks ($50{+}$ layers) by ensuring gradients in the residual stream don't explode/vanish.
At our depth ($L=10$), this is less critical.
However, pre-norm APN could help for:
\begin{itemize}
\item Multi-head APN (where instability is the main problem).
\item Deeper stacks ($L \geq 20$).
\end{itemize}
\end{enumerate}

\subsection{Would Switching to Pre-Norm Help?}

Possibly, with caveats. A pre-norm APN would look like:
\begin{equation}
\mathbf{h} \leftarrow \mathbf{h} + \text{APN}(\text{LN}(\mathbf{h}))
\end{equation}
with the internal \texttt{o\_norm} removed. The risk is that normalizing the input keys reduces APN's effective associative capacity (fewer distinguishable keys in the normalized space).

This is testable and we plan to investigate it in future work, particularly for multi-head configurations where stability is the primary concern. The hypothesis is that pre-norm + constrained $\eta$ (via softplus) together might enable stable multi-head APN.

% ============================================================================
\section{Summary of Architectural Choices}

\begin{table}[h]
\centering
\caption{Normalization placement across all models in our experiments.}
\begin{tabular}{lll}
\toprule
\textbf{Model} & \textbf{Norm Type} & \textbf{Formula} \\
\midrule
Transformer & Pre-norm & $\mathbf{h} + \text{MHA}(\text{LN}(\mathbf{h}))$; $\mathbf{h} + \text{FFN}(\text{LN}(\mathbf{h}))$ \\
APN (no FFN) & Branch-norm & $\mathbf{h} + \text{LN}(\text{APN\_output}(\mathbf{h}))$ \\
APN + FFN & Branch + Pre & $\mathbf{h} + \text{LN}(\text{APN\_output}(\mathbf{h}))$; $\mathbf{h} + \text{FFN}(\text{LN}(\mathbf{h}))$ \\
GFR (no FFN) & Branch-norm & $\mathbf{h} + \text{LN}(\text{GFR\_output}(\mathbf{h}))$ \\
DeltaNet & Pre-norm & $\mathbf{h} + \text{DeltaNet}(\text{LN}(\mathbf{h}))$ \\
\midrule
All models & Final norm & $W_{\text{cls}} \cdot \text{LN}(\mathbf{h}_T)$ \\
\bottomrule
\end{tabular}
\end{table}

\subsection{Future Experiments Motivated by These Questions}

\begin{enumerate}
\item \textbf{Constrained $\eta$ via softplus}: Should fix multi-head divergence while preserving the anti-Hebbian correction capability (which comes from the sign of $\mathbf{v} - S^\top\mathbf{k}$, not from the sign of $\eta$).
\item \textbf{Pre-norm APN}: Test whether $\mathbf{h} + \text{APN}(\text{LN}(\mathbf{h}))$ improves stability at the cost of key distinguishability.
\item \textbf{Oja-like feature normalization}: Apply Oja's normalization to the \emph{representations} before they enter the delta-rule memory, as a complementary regularizer---this could combine the best of both approaches.
\end{enumerate}

\end{document}
